% !Mode:: "TeX:UTF-8"
% !TEX program  = xelatex

\section{补充实验配置与复现说明}

\subsection{动机性预分析的详细统计}

本节汇总引言中动机性实验对应的详细统计。

\paragraph{TSPLIB 分析：不存在统一支配所有实例的方法。}
表 \ref{tab:motivation-tsplib-overall} 与表 \ref{tab:motivation-tsplib-by-size} 
汇总了 61 个 TSPLIB 实例上的初始化方法比较。尽管 LEHD 在整体上胜出次数最多，
但方法之间不存在完全的支配关系，按规模分类后也未观察到规模偏好。

\begin{table}[htb]
  \centering
  \caption{61 个 TSPLIB 实例上的初始化方法预分析}
  \label{tab:motivation-tsplib-overall}
  \small
  \begin{tabular}{lcccc}
    \toprule
    方法 & 实例数 & win(pct) & tie & 平均 gap \\
    \midrule
    lehd & 61 & 43 (70.49\%) & 5 & 2.67 \\
    elg  & 61 & 7 (11.48\%)  & 2 & 4.93 \\
    icam & 61 & 4 (6.56\%)   & 2 & 5.29 \\
    omni & 61 & 2 (3.27\%)   & 1 & 8.18 \\
    dact & 61 & 0            & 0 & 25.18 \\
    \bottomrule
  \end{tabular}
\end{table}

\begin{table}[htb]
  \centering
  \caption{61 个 TSPLIB 实例上按规模分类的最佳方法分布}
  \label{tab:motivation-tsplib-by-size}
  \small
  \begin{tabular}{lccccc}
    \toprule
    规模区间 $n$ & 实例数 & lehd & elg & icam & omni \\
    \midrule
    $\leq 100$     & 12 & 10 (83.3\%) & 3 (25.0\%) & 3 (25.0\%) & 1 (8.3\%) \\
    101--200       & 17 & 12 (70.6\%) & 3 (17.6\%) & 0 (0.0\%)  & 2 (11.8\%) \\
    201--500       & 13 & 9 (69.2\%)  & 1 (7.7\%)  & 3 (23.1\%) & 0 (0.0\%)  \\
    501--1000      & 6  & 5 (83.3\%)  & 1 (16.7\%) & 0 (0.0\%)  & 0 (0.0\%)  \\
    1001--2000     & 13 & 12 (92.3\%) & 1 (7.7\%)  & 0 (0.0\%)  & 0 (0.0\%)  \\
    \bottomrule
  \end{tabular}
\end{table}

\paragraph{合成 TSP 数据集分析：分布变化改变领先方法。}
表 \ref{tab:motivation-synth-tsp} 给出了三组合成 TSP 预分析中的胜出计数。
不同分布下的领先方法并不固定，说明方法对实例具有偏好性。

\begin{table}[htb]
  \centering
  \caption{合成 TSP 预分析中的胜出计数}
  \label{tab:motivation-synth-tsp}
  \small
  \begin{tabular}{llcc}
    \toprule
    数据设置 & 方法 & wins(aug) & wins(no\_aug) \\
    \midrule
    tsp2000-uniform    & omni & 0   & 0  \\
    tsp2000-uniform    & lehd & 4   & 32 \\
    tsp2000-uniform    & elg  & 0   & 0  \\
    tsp2000-uniform    & icam & 124 & 96 \\
    \midrule
    tsp1000-explosion  & omni & 0   & 0  \\
    tsp1000-explosion  & lehd & 100 & 100 \\
    tsp1000-explosion  & elg  & 0   & 0  \\
    tsp1000-explosion  & icam & 0   & 0  \\
    \midrule
    tsp1000-rotation   & omni & 0   & 0  \\
    tsp1000-rotation   & lehd & 99  & 99 \\
    tsp1000-rotation   & elg  & 0   & 0  \\
    tsp1000-rotation   & icam & 1   & 1  \\
    \bottomrule
  \end{tabular}
\end{table}

\paragraph{CVRPLIB 分析：选择需求同样存在于另一类问题。}
表 \ref{tab:motivation-cvrp-overall} 与表 \ref{tab:motivation-cvrp-by-size} 
汇总了 100 个 CVRPLIB 实例上的结果。
ICAM 与 ReLD 在整体上几乎相同，但不同规模段的领先方法又发生变化。这说明不同规模段下的领先方法并不固定。

\begin{table}[htb]
  \centering
  \caption{100 个 CVRPLIB 实例上的初始化方法预分析}
  \label{tab:motivation-cvrp-overall}
  \small
  \begin{tabular}{lccc}
    \toprule
    方法 & win(pct) & tie & 平均 gap \\
    \midrule
    icam & 46 (46\%) & 0 & 4.34 \\
    reld & 46 (46\%) & 0 & 6.09 \\
    omni & 8 (8\%)   & 0 & 4.02 \\
    \bottomrule
  \end{tabular}
\end{table}

\begin{table}[htb]
  \centering
  \caption{100 个 CVRPLIB 实例上按规模分类的最佳方法分布}
  \label{tab:motivation-cvrp-by-size}
  \small
  \begin{tabular}{lcccc}
    \toprule
    规模区间 $n$ & 实例数 & icam & reld & omni \\
    \midrule
    $\leq 200$ & 22 & 5 (22.7\%)  & 14 (63.6\%) & 3 (13.6\%) \\
    201--500   & 46 & 23 (50.0\%) & 19 (41.3\%) & 4 (8.7\%)  \\
    501--800   & 22 & 13 (59.1\%) & 8 (36.4\%)  & 1 (4.5\%)  \\
    $> 800$    & 10 & 5 (50.0\%)  & 5 (50.0\%)  & 0 (0.0\%)  \\
    \bottomrule
  \end{tabular}
\end{table}

\paragraph{组合实验：初始化与迭代确实存在耦合。}
表 \ref{tab:motivation-composition} 给出了 TSPLIB 上四组“初始化+迭代”组合的直接比较。
结果显示，替换迭代器或初始化器，都能在一部分实例上带来提升。

\begin{table}[htb]
  \centering
  \caption{TSPLIB 上初始化与迭代组合的直接比较}
  \label{tab:motivation-composition}
  \small
  \begin{tabular}{lccc}
    \toprule
    组合对比 & A\_win & tie & B\_win \\
    \midrule
    glop+glop vs. glop+lehd & 9  & 2  & 50 \\
    lehd+lehd vs. lehd+glop & 37 & 10 & 14 \\
    glop+glop vs. lehd+glop & 14 & 1  & 46 \\
    lehd+lehd vs. glop+lehd & 32 & 13 & 16 \\
    \bottomrule
  \end{tabular}
\end{table}

\subsection{训练超参数}

表 \ref{tab:default-hparams} 汇总了本文主要上下文老虎机实现中 Neural-LinUCB 的超参数。

\begin{table}[htb]
  \centering
  \caption{Neural-LinUCB 超参数}
  \label{tab:default-hparams}
  \begin{tabular}{lc}
    \toprule
    参数 & 数值 \\
    \midrule
    hidden\_dim & 64 \\
    arm\_embed\_dim & 16 \\
    $\alpha$ & 1.0 \\
    reg & 1.0 \\
    lr & $1\times 10^{-3}$ \\
    train\_every & 100 \\
    representation\_steps & 10 \\
    representation\_buffer\_size & 5000 \\
    representation\_batch\_size & 512 \\
    train\_epsilon & 0.10 \\
    initial\_pulls & 2 \\
    \bottomrule
  \end{tabular}
\end{table}

\subsection{计算资源与训练时间}

所有实验均在单张 NVIDIA GeForce RTX 3090 GPU（24 GB 显存）上完成。
Neural-LinUCB 在合成混合数据集上的完整训练时间约为 5 小时。